\section{Worlds Apart}

\begin{quotex}
Opportunities for interpenetration have never been greater---and yet our minds are worlds apart.


\flright{\textsc{Owen Barfield}}
\end{quotex}


\begin{quotex}
It is a platitude to say that the men of our day no longer understand each other. This is nothing new since the Tower of Babel. But God had left a certain relation between men through the intelligence. Now it is this bridge which our day has just broken down. Men no longer understand each other: on the level of the peasant this is not evident; on the bourgeois level is an inconvenience, on the intellectual level it is tragic, because, for the intellectual there is no other genuine reason for living than that of communication, in order to understand the world. Today this communication has become practically impossible. In order to understand each other we need a minimum of ideas which are common and valid for everyone, of prejudices and values which are the same for all---and most often unconscious.


\flright{\textsc{Jacques Ellul}, \emph{The Presence of the Kingdom}}
\end{quotex}


After thousands of years, millions of books, billions of words, trillions of thoughts, mankind is still no closer to a common understanding of reality. Along with Solovyov, we can trace this to mechanical thinking. Discourse used to be organic, it was verbal and participatory, it led to a change of being.

The written word was the beginning of the change, a change to mechanical thinking. As technics became the media of communications, discourse came under the domination of technicians. Since technics are mechanical, thinking came to be limited to the mechanical. Thought is rationalized, and anything that does not fit in, is discounted as mere feeling, subjective idea, superstition, or worse.

A similar situation arose in the religious sphere. True faith is organic, involving the whole being. More than mere belief, faith includes the will and feeling. But when the discourse on faith was taken over by the half-educated mechanical mind, it could only understand faith as belief. They then denied that the object of faith could not be comprehended or directly known. Thus, belief became institutionalized.

The rationalists and the fideists each claimed to be the arbiters of truth, and claimed sole competence to define the limits of knowledge. For two thousand years they have been battling each other in various ways, while true seekers of knowledge are marginalized and relegated to the fringe. The best they can hope for is to be left alone to pursue their studies.

Organic knowledge is possible only to the integral personality, that is, a person with the thinking, willing and feeling functions acting harmoniously, each with its place in knowledge. For each function there is a corresponding object and subject.

\begin{table}[h]\centering\small
\begin{tabular}{ccc}\toprule
\textbf{Function} & \textbf{Object} & \textbf{Subject} \\\midrule
Willing & Goodness & Sprit \\ Thinking & Truth & Intellect \\ Feeling & Beauty & Soul \\\bottomrule
\end{tabular}
\end{table}

In organic thinking all the functions work together in knowing. The world is will and representation, and is felt as well as known. Mechanical thinking is limited to the intellect, and hence incapable of knowing the Good and the Beautiful, which it can only see as subjective and arbitrary.

Reference:

\emph{Lectures on Divine Humanity}, by Vladimir Solovyov



\flrightit{Posted on 2010-04-12 by Cologero}

